\IEEEPARstart{A}{utonomous} vehicles plan around forecasts of what nearby road users will do. A forecast is only as
useful as the statement it makes about its own reliability: a planner that knows a predicted region contains the true
future with probability at least $1-\alpha$ can keep clear of that region and of nothing larger. Conformal prediction
(CP) supplies such a statement for any black-box predictor, with a finite-sample guarantee that requires nothing beyond
exchangeability of calibration and test data \cite{vovk2005algorithmic,angelopoulos2023gentle}, and it has been adopted
for trajectory prediction and safe planning
\cite{lindemann2023safe,sun2024copula,huang2025cuqds}.

Exchangeability is exactly what deployment breaks. A forecaster is calibrated on the data at hand---one set of cities,
one sensor suite, one annotation pipeline---and then driven somewhere else. That accuracy drops when a trajectory
predictor moves between datasets is documented \cite{feng2024unitraj}, and that predictive uncertainty degrades under
dataset shift is well known for classifiers \cite{ovadia2019can}. What is not known is what happens to the
\emph{guarantee itself} when a conformally calibrated forecaster is moved, without any new labels, to a different real
driving dataset: whether the stated $1-\alpha$ survives, how badly it fails, and whether the standard remedies for shift
repair it.

The conformal literature offers several remedies, each resting on an assumption. Weighted CP restores validity under
covariate shift when the likelihood ratio is known or can be estimated \cite{tibshirani2019covariate}; more general
results bound the coverage loss by a distance between calibration and test distributions
\cite{barber2023beyond}; online methods track shift by feeding realised errors back into the threshold
\cite{gibbs2021adaptive,angelopoulos2023pid}, an idea carried to trajectory prediction \cite{huang2025cuqds} and to
robotic settings more broadly \cite{tumu2026adaptnc}; and a generative-prior approach has been proposed for planning
under environment shift \cite{rahaman2026shift}. What these do not establish is the \emph{static, label-free,
dataset-to-dataset} setting in which a forecaster is frozen, calibrated once, and deployed: how large the failure is on
real data, whether the covariate-shift remedy suffices, and how many labelled target scenes it takes to detect or repair
the failure.

\paragraph*{Contributions}
\todo{Finalise after results. Planned, each to be kept ONLY if the pre-registered test supports it: (i) a measurement protocol and the first measurement of zero-label conformal coverage transfer between real driving datasets; (ii) a diagnosis --- covariate reweighting does not repair the loss --- with a proposition explaining why label-free repair is limited; (iii) the solution(s) that worked: model-uncertainty-normalised scores (H4), a label-efficient coverage audit with a closed-form label budget (H5b), a label-free shift monitor (H6); (iv) released code, pre-registration and experiment log.}
